\documentclass[../../main.tex]{subfiles}

\begin{document}

\section{$L^{p}$ 中的弱序列收敛}\label{section:ch08-弱序列收敛}

关于实数的 Bolzano-Weierstrass 定理断言每个有界实数序列有一个收敛的子序列. 该性质可立即推广到每个欧氏空间 $\mathbf{R}^{n}$ 的有界序列. 该性质在无穷维赋范线性空间中不成立\footnote{13.3 节中证明的 Riesz 定理告诉我们, 在每个无穷维赋范线性空间 $X$ 中, 存在有界序列, 它没有在 $X$ 中收敛的子序列.}. 特别地, 以下例子表明对于 $1\leqslant p\leqslant\infty$, 有 $L^{p}[0,1]$ 中的有界序列不存在在 $L^{p}[0,1]$ 中收敛的子序列. 在以下例子中定义的函数称为 Rademacher 函数.

\begin{example}\label{example:ch08-Rademacher}
对于 $I=[0,1]$ 以及自然数 $k$ 和 $n$, 考虑定义在 $I$ 上的阶梯函数 $f_{n}$,
\begin{align*}
f_{n}(x)=(-1)^{k},\quad k/2^{n}\leqslant x<(k+1)/2^{n},\quad\text{其中 }0\leqslant k<2^{n}-1.
\end{align*}
固定 $1\leqslant p\leqslant\infty$. 则 $\{f_{n}\}$ 是 $L^{p}(I)$ 中的有界序列: 的确, 对每个指标 $n$, $\|f_{n}\|_{p}\leqslant 1$. 另一方面, 由于对 $n\neq m$, $|f_{n}-f_{m}|$ 在测度为 $1/2$ 的集合上取值 $2$, $\|f_{n}-f_{m}\|_{p}\geqslant(2)^{1-1/p}$. 因此 $\{f_{n}\}$ 没有在 $L^{p}(I)$ 中是 Cauchy 的子序列, 且因此没有在 $L^{p}(I)$ 中收敛的子序列. 我们也注意到没有子序列在 $I$ 上几乎处处逐点收敛, 这是由于对 $1\leqslant p<\infty$, 若存在这样的子序列, 根据有界收敛定理, 它将在 $L^{p}(I)$ 中收敛.
\end{example}

\begin{definition}[弱收敛]\label{definition:ch08-弱收敛}
令 $X$ 为赋范线性空间. $X$ 中的序列 $\{f_{n}\}$ 称为在 $X$ 中弱收敛于 $X$ 中的 $f$, 若
\begin{align*}
\text{对所有 }T\in X^{*},\quad\lim_{n\to\infty}T(f_{n})=T(f).
\end{align*}
我们记作:
\begin{align*}
\text{在 }X\text{ 中},\quad\{f_{n}\}\rightharpoonup f,
\end{align*}
意味着 $f$ 和每个 $f_{n}$ 属于 $X$, 且 $\{f_{n}\}$ 在 $X$ 中弱收敛于 $f$.
\end{definition}

我们继续用在 $X$ 中 $\{f_{n}\}\to f$ 表示 $\lim_{n\to\infty}\|f_{n}-f\|=0$, 为区别这种模式的收敛与弱收敛, 常称这种模式的收敛为在 $X$ 中强收敛. 由于对所有 $T\in X^{*}$,
\begin{align*}
|T(f_{n})-T(f)|=|T(f_{n}-f)|\leqslant\|T\|_{*}\cdot\|f_{n}-f\|.
\end{align*}
若一个序列强收敛, 则它弱收敛. 反过来不成立.

\begin{proposition}\label{proposition:ch08-Lp弱收敛刻画}
令 $E$ 为可测集, $1\leqslant p<\infty$, 而 $q$ 为 $p$ 的共轭. 则在 $L^{p}(E)$ 中, $\{f_{n}\}\rightharpoonup f$ 当且仅当对所有 $g\in L^{q}(E)$,
\begin{align*}
\lim_{n\to\infty}\int_{E}g\cdot f_{n}=\int_{E}g\cdot f.
\end{align*}
\end{proposition}

\begin{proof}
Riesz 表示定理告诉我们 $L^{p}(E)$ 上的每个有界线性泛函由 $L^{q}(E)$ 中的某个函数的积分给出.
\end{proof}

对于可测集 $E$ 与 $1\leqslant p<\infty$, $L^{p}(E)$ 中的序列至多只能弱收敛于 $L^{p}(E)$ 中的一个函数. 事实上, 假定 $\{f_{n}\}$ 在 $L^{p}(E)$ 中弱收敛于 $f_{1}$ 和 $f_{2}$. 考虑 $f_{1}-f_{2}$ 的共轭函数 $(f_{1}-f_{2})^{*}$. 则
\begin{align*}
\int_{E}(f_{1}-f_{2})^{*}\cdot f_{2}=\lim_{n\to\infty}\int_{E}(f_{1}-f_{2})^{*}\cdot g_{n}=\int_{E}(f_{1}-f_{2})^{*}\cdot f_{1}.
\end{align*}
因此
\begin{align*}
\|f_{1}-f_{2}\|_{p}=\int_{E}(f_{1}-f_{2})^{*}(f_{1}-f_{2})=0.
\end{align*}
于是 $f_{1}=f_{2}$, 从而弱收敛的极限是唯一的.

\begin{theorem}\label{theorem:ch08-弱收敛有界}
令 $E$ 为可测集而 $1\leqslant p<\infty$. 假定在 $L^{p}(E)$ 中 $\{f_{n}\}\rightharpoonup f$. 则
\begin{align}
\{f_{n}\}\text{ 在 }L^{p}(E)\text{ 中是有界的且 }\|f\|_{p}\leqslant\liminf\|f_{n}\|_{p}.\label{align:ch08-弱收敛有界}
\end{align}
\end{theorem}

\begin{proof}
令 $q$ 为 $p$ 的共轭而 $f^{*}$ 为 $f$ 的共轭函数. 我们首先证明\eqref{align:ch08-弱收敛有界}右边的不等式. 我们从 H\"older 不等式推出, 对所有 $n$,
\begin{align*}
\int_{E}f^{*}\cdot f_{n}\leqslant\|f^{*}\|_{q}\cdot\|f_{n}\|_{p}=\|f_{n}\|_{p}.
\end{align*}
由于 $\{f_{n}\}$ 弱收敛于 $f$ 而 $f^{*}$ 属于 $L^{q}(E)$,
\begin{align*}
\|f\|_{p}=\int_{E}f^{*}\cdot f=\lim_{n\to\infty}\int_{E}f^{*}\cdot f_{n}\leqslant\liminf\|f_{n}\|_{p}.
\end{align*}
我们用反证法证明 $\{f_{n}\}$ 在 $L^{p}(E)$ 中是有界的. 假设 $\{\|f_{n}\|_{p}\}$ 无界. 不失一般性 (见习题 18), 通过取一个子序列的数乘, 我们假定对所有 $n$,
\begin{align}
\|f_{n}\|_{p}=n\cdot 3^{n}.\label{align:ch08-无界假设}
\end{align}
我们归纳地选取实数序列 $\{\varepsilon_{k}\}$ 使得对每个 $k$, $\varepsilon_{k}=\pm 1/3$. 定义 $\varepsilon_{1}=1/3$. 若 $n$ 是一个自然数使得 $\varepsilon_{1},\cdots,\varepsilon_{n}$ 已被定义, 定义
\begin{align*}
\varepsilon_{n+1}=1/3^{n+1},\quad\text{若 }\int_{E}\left[\sum_{k=1}^{n}\varepsilon_{k}(f_{k})^{*}\right]\cdot f_{n+1}\geqslant 0,
\end{align*}
而 $\varepsilon_{n+1}=-1/3^{n+1}$, 若上述积分是负的. 因此, 根据\eqref{align:ch08-无界假设}和共轭函数的定义, 对所有 $n$,
\begin{align}
\left|\int_{E}\left[\sum_{k=1}^{n}\varepsilon_{k}(f_{k})^{*}\right]\cdot f_{n}\right|\geqslant 1/3^{n}\|f_{n}\|_{p}=n\text{ 且 }\|\varepsilon_{n}\cdot(f_{n})^{*}\|_{q}=1/3^{n}.\label{align:ch08-弱收敛矛盾}
\end{align}
由于对所有 $k$, $\|\varepsilon_{k}\cdot(f_{k})^{*}\|_{q}=1/3^{k}$, 级数 $\sum_{k=1}^{\infty}\varepsilon_{k}\cdot(f_{k})^{*}$ 的部分和序列是 $L^{q}(E)$ 中的 Cauchy 序列. Riesz-Fischer 定理告诉我们 $L^{q}(E)$ 是完备的. 定义函数 $g\in L^{q}(E)$ 为
\begin{align*}
g=\sum_{k=1}^{\infty}\varepsilon_{k}\cdot(f_{k})^{*}.
\end{align*}
固定自然数 $n$. 我们从三角不等式\eqref{align:ch08-弱收敛矛盾}和 H\"older 不等式推出
\begin{align*}
\left|\int_{E}g\cdot f_{n}\right|&=\left|\int_{E}\left[\sum_{k=1}^{\infty}\varepsilon_{k}\cdot(f_{k})^{*}\right]\cdot f_{n}\right|\\
&\geqslant\left|\int_{E}\left[\sum_{k=1}^{n}\varepsilon_{k}\cdot(f_{k})^{*}\right]\cdot f_{n}\right|-\left|\int_{E}\left[\sum_{k=n+1}^{\infty}\varepsilon_{k}\cdot(f_{k})^{*}\right]\cdot f_{n}\right|\\
&\geqslant n-\left|\int_{E}\left[\sum_{k=n+1}^{\infty}\varepsilon_{k}\cdot(f_{k})^{*}\right]\cdot f_{n}\right|\\
&\geqslant n-\left[\sum_{k=n+1}^{\infty}1/3^{k}\right]\cdot\|f_{n}\|_{p}\\
&=n-1/3^{n}\cdot 1/2\cdot\|f_{n}\|_{p}=n/2.
\end{align*}
这是一个矛盾, 因为序列 $\{f_{n}\}$ 在 $L^{p}(E)$ 中弱收敛而 $g$ 属于 $L^{q}(E)$, 实数序列 $\{\int_{E}g\cdot f_{n}\}$ 收敛, 因而是有界的. 因此 $\{f_{n}\}$ 在 $L^{p}$ 中是有界的.
\end{proof}

\begin{corollary}\label{corollary:ch08-弱收敛乘积}
令 $E$ 为可测集, $1\leqslant p<\infty$, 而 $q$ 为 $p$ 的共轭. 假定 $\{f_{n}\}$ 在 $L^{p}(E)$ 中弱收敛于 $f$, 而 $\{g_{n}\}$ 在 $L^{q}(E)$ 中强收敛于 $g$. 则
\begin{align}
\lim_{n\to\infty}\int_{E}g_{n}\cdot f_{n}=\int_{E}g\cdot f.\label{align:ch08-弱收敛乘积}
\end{align}
\end{corollary}

\begin{proof}
对每个指标 $n$,
\begin{align*}
\int_{E}g_{n}\cdot f_{n}-\int_{E}g\cdot f=\int_{E}[g_{n}-g]\cdot f_{n}+\int_{E}g\cdot f_{n}-\int_{E}g\cdot f.
\end{align*}
根据前一个定理, 存在常数 $C\geqslant 0$ 使得对所有 $n$,
\begin{align*}
\|f_{n}\|_{p}\leqslant C.
\end{align*}
因此, 根据 H\"older 不等式, 对所有 $n$,
\begin{align*}
\left|\int_{E}g_{n}\cdot f_{n}-\int_{E}g\cdot f\right|\leqslant C\cdot\|g_{n}-g\|_{q}+\left|\int_{E}g\cdot f_{n}-\int_{E}g\cdot f\right|.
\end{align*}
从这些不等式和事实
\begin{align*}
\lim_{n\to\infty}\|g_{n}-g\|_{q}=0\text{ 且 }\lim_{n\to\infty}\int_{E}g\cdot f_{n}=\int_{E}g\cdot f
\end{align*}
得出\eqref{align:ch08-弱收敛乘积}成立.
\end{proof}

提到线性空间 $X$ 的子集 $\mathcal{S}$ 的线性扩张, 我们指的是由 $\mathcal{S}$ 中函数的所有线性组合组成的线性空间, 即形如
\begin{align*}
f=\sum_{k=1}^{n}\alpha_{k}\cdot f_{k}
\end{align*}
的函数的线性空间, 其中每个 $\alpha_{k}$ 是实数, 而每个 $f_{k}$ 属于 $\mathcal{S}$.

\begin{proposition}\label{proposition:ch08-稠密集弱收敛}
令 $E$ 为可测集, $1\leqslant p<\infty$, 而 $q$ 为 $p$ 的共轭. 假设 $\mathcal{F}$ 是 $L^{q}(E)$ 的一个子集, 其线性扩张在 $L^{q}(E)$ 中稠密. 令 $\{f_{n}\}$ 为 $L^{p}(E)$ 中的有界序列而 $f$ 属于 $L^{p}(E)$. 则在 $L^{p}(E)$ 中 $\{f_{n}\}\rightharpoonup f$ 当且仅当对所有 $g\in\mathcal{F}$,
\begin{align}
\lim_{n\to\infty}\int_{E}f_{n}\cdot g=\int_{E}f\cdot g.\label{align:ch08-稠密集弱收敛}
\end{align}
\end{proposition}

\begin{proof}
命题 6 刻画了 $L^{q}(E)$ 中的弱收敛. 假设\eqref{align:ch08-稠密集弱收敛}成立. 为验证弱收敛性, 令 $g_{0}$ 属于 $L^{q}(E)$. 我们证明 $\lim_{n\to\infty}\int_{E}f_{n}\cdot g_{0}=\int_{E}f\cdot g_{0}$. 令 $\varepsilon>0$. 我们必须找到自然数 $N$ 使得对 $n\geqslant N$,
\begin{align}
\left|\int_{E}f_{n}\cdot g_{0}-\int_{E}f\cdot g_{0}\right|<\varepsilon.\label{align:ch08-稠密集估计}
\end{align}
观察到对任何 $g\in L^{q}(E)$ 和自然数 $n$,
\begin{align*}
\int_{E}f_{n}\cdot g_{0}-\int_{E}f\cdot g_{0}=\int_{E}(f_{n}-f)\cdot(g_{0}-g)+\int_{E}(f_{n}-f)\cdot g.
\end{align*}
因此, 根据 H\"older 不等式,
\begin{align*}
\left|\int_{E}f_{n}\cdot g_{0}-\int_{E}f\cdot g_{0}\right|\leqslant\|f_{n}-f\|_{p}\cdot\|g-g_{0}\|_{q}+\left|\int_{E}f_{n}\cdot g-\int_{E}f\cdot g\right|.
\end{align*}
由于 $\{f_{n}\}$ 在 $L^{p}(E)$ 中有界, 而 $\mathcal{F}$ 的线性扩张在 $L^{q}(E)$ 中稠密, 该线性扩张中存在函数 $g$ 使得对所有 $n$, $\|f_{n}-f\|_{p}\cdot\|g-g_{0}\|_{q}<\varepsilon/2$.

但 $g$ 是 $\mathcal{F}$ 中函数的线性组合, 因此我们从\eqref{align:ch08-稠密集弱收敛}推出
\begin{align*}
\lim_{n\to\infty}\int_{E}f_{n}\cdot g=\int_{E}f\cdot g.
\end{align*}
因此存在自然数 $N$ 使得若 $n\geqslant N$,
\begin{align*}
\left|\int_{E}f_{n}\cdot g-\int_{E}f\cdot g\right|<\varepsilon/2.
\end{align*}
根据前一个估计显然\eqref{align:ch08-稠密集估计}对 $N$ 的这个选择成立.
\end{proof}

根据前一章的命题 9, 对于 $1\leqslant q\leqslant\infty$, $L^{p}(E)$ 中的简单函数在 $L^{p}(E)$ 中稠密, 且若 $q<\infty$, 这些函数有有限支撑. 此外, 前一章的命题 10 告诉我们对于闭有界区间 $[a,b]$ 和 $1<q<\infty$, 阶梯函数在 $L^{q}[a,b]$ 中稠密. 因此弱连续性的以下两个刻画从前一个命题得到.

\begin{theorem}\label{theorem:ch08-特征函数判别}
令 $E$ 为可测集而 $1\leqslant p<\infty$. 假定 $\{f_{n}\}$ 是 $L^{p}(E)$ 中的有界序列而 $f$ 属于 $L^{p}(E)$. 则在 $L^{p}(E)$ 中 $\{f_{n}\}\rightharpoonup f$ 当且仅当对 $E$ 的每个可测子集 $A$,
\begin{align}
\lim_{n\to\infty}\int_{A}f_{n}=\int_{A}f.\label{align:ch08-特征函数判别}
\end{align}
若 $p>1$, 仅须考虑有限测度集 $A$.
\end{theorem}

\begin{theorem}\label{theorem:ch08-区间积分判别}
令 $[a,b]$ 为闭有界区间而 $1<p<\infty$. 假定 $\{f_{n}\}$ 是 $L^{p}[a,b]$ 中的有界序列而 $f$ 属于 $L^{p}[a,b]$. 则在 $L^{p}[a,b]$ 中, $\{f_{n}\}\rightharpoonup f$ 当且仅当
\begin{align}
\text{对所有 }[a,b]\text{ 中的 }x,\quad\lim_{n\to\infty}\left[\int_{a}^{x}f_{n}\right]=\int_{a}^{x}f.\label{align:ch08-区间积分判别}
\end{align}
\end{theorem}

由于阶梯函数在 $L^{\infty}[a,b]$ 中不稠密 (见习题 44), 当 $p=1$ 时定理 11 不成立.

\begin{example}[Riemann-Lebesgue 引理]\label{example:ch08-Riemann-Lebesgue}
令 $I=[-\pi,\pi]$ 而 $1<p<\infty$. 对每个自然数 $n$, 对 $I$ 中的 $x$ 定义 $f_{n}(x)=\sin nx$. 我们宣称 $\{f_{n}\}$ 在 $L^{p}(I)$ 中弱收敛于 $f\equiv 0$. 事实上, 观察到对每个 $n$, 在 $I$ 上 $|f_{n}|\leqslant 1$ 且对所有 $x\in I$,
\begin{align*}
\lim_{n\to\infty}\int_{-\pi}^{x}\sin nt\,dt=\lim_{n\to\infty}-\frac{1}{n}[\cos nx-(-1)^{n}]=0.
\end{align*}
因此, 对 $1<p<\infty$, 我们从定理 11 推出 $\{f_{n}\}$ 在 $L^{p}(I)$ 中弱收敛于 $f$. 对于 $p=1$ 的情形, 弱收敛从定理 10、简单函数在 $L^{\infty}$ 中的稠密性以及第 2 章的定理 12 得出. 观察到对每个 $n$,
\begin{align*}
\int_{-\pi}^{\pi}|\sin nt|^{2}\,dt=\int_{-\pi}^{\pi}\sin^{2}nt\,dt=\pi.
\end{align*}
因此 $\{f_{n}\}$ 没有在 $L^{2}(I)$ 中强收敛到 $f$ 的子序列. 类似的估计表明 $\{f_{n}\}$ 在任何 $L^{p}(I)$ ($1\leqslant p<\infty$) 中没有强收敛的子序列.
\end{example}

\begin{example}\label{example:ch08-逐点不弱收敛}
对于自然数 $n$, 在 $[0,1]$ 上定义 $f_{n}=n\cdot\chi_{(0,1/n)}$. 在 $[0,1]$ 上定义 $f\equiv 0$. 则 $\{f_{n}\}$ 是 $L^{1}[0,1]$ 中的单位函数序列. 它在 $[0,1]$ 上逐点收敛于 $f$. 但在 $L^{1}[0,1]$ 中 $\{f_{n}\}$ 不弱收敛于 $f$, 这是由于取 $g=\chi_{[0,1]}\in L^{\infty}[0,1]$,
\begin{align*}
\lim_{n\to\infty}\int_{0}^{1}g\cdot f_{n}=\lim_{n\to\infty}\int_{0}^{1}f_{n}=1,\quad\text{而}\int_{0}^{1}g\cdot f=\int_{0}^{1}f=0.
\end{align*}
\end{example}

\begin{example}\label{example:ch08-平移帐篷}
定义 $\mathbf{R}$ 上在 $(-1,1)$ 外消失、在区间 $[-1,0]$ 和 $[0,1]$ 上线性且在 $x=0$ 取值为 $1$ 的帐篷函数 $f_{0}$. 对每个自然数 $n$, 定义 $f_{n}(x)=f_{0}(x-n)$ 且在 $\mathbf{R}$ 上令 $f\equiv 0$. 则在 $\mathbf{R}$ 上逐点地 $f_{n}\to f$. 令 $1\leqslant p<\infty$. 序列 $\{f_{n}\}$ 在 $L^{p}(\mathbf{R})$ 中是有界的. 我们把以下结论的证明留作练习: 用测度的连续性证明对有限测度的集合 $A$,
\begin{align}
\lim_{n\to\infty}\int_{A}f_{n}=\int_{A}f.\label{align:ch08-平移帐篷}
\end{align}
从而从定理 10 推出, 对 $1<p<\infty$, 在 $L^{p}(\mathbf{R})$ 中 $\{f_{n}\}\rightharpoonup f$. 但在 $L^{1}(\mathbf{R})$ 中 $\{f_{n}\}$ 不弱收敛于 $f$, 这是由于对 $\mathbf{R}$ 上的函数 $g\equiv 1$, $g$ 属于 $L^{\infty}(\mathbf{R})$, 而 $\{\int_{\mathbf{R}}f_{n}\}$ 不收敛于 $\int_{\mathbf{R}}f$.
\end{example}

前两个例子展示了 $L^{1}(E)$ 中逐点收敛于 $L^{1}(E)$ 中的函数的有界序列, 然而在 $L^{1}(E)$ 中不弱收敛. 若 $1<p<\infty$, 这在 $L^{p}(E)$ 中不可能发生.

\begin{theorem}\label{theorem:ch08-逐点有界弱收敛}
令 $E$ 为可测集而 $1<p<\infty$. 假定 $\{f_{n}\}$ 是 $L^{p}(E)$ 中在 $E$ 上 a.e. 逐点收敛到 $f$ 的一个有界序列. 则在 $L^{p}(E)$ 中 $\{f_{n}\}\rightharpoonup f$.
\end{theorem}

\begin{proof}
我们将 Fatou 引理用到序列 $\{|f_{n}|^{p}\}$, 推出 $f$ 属于 $L^{p}(E)$. 定理 10 告诉我们验证弱序列收敛必要且充分的是证明对 $E$ 的每个有限测度的可测子集 $A$,
\begin{align}
\lim_{n\to\infty}\int_{A}f_{n}=\int_{A}f.\label{align:ch08-逐点有界判别}
\end{align}
令 $A$ 为这样的一个子集. 根据前一章的推论 2, 由于序列 $\{f_{n}\}$ 在 $L^{p}(E)$ 中是有界的, 它在 $E$ 上是一致可积的. 但 $m(A)<\infty$. 因此根据 Vitali 收敛定理, \eqref{align:ch08-逐点有界判别}成立.
\end{proof}

\begin{theorem}[Radon-Riesz 定理]\label{theorem:ch08-Radon-Riesz}
令 $E$ 为可测集而 $1<p<\infty$. 假定在 $L^{p}(E)$ 中 $\{f_{n}\}\rightharpoonup f$. 则
\begin{align*}
\text{在 }L^{p}(E)\text{ 中 }\{f_{n}\}\to f\text{ 当且仅当 }\lim_{n\to\infty}\|f_{n}\|_{p}=\|f\|_{p}.
\end{align*}
\end{theorem}

\begin{proof}
在任何赋予范数 $\|\cdot\|$ 的线性空间 $X$ 中, 我们总是有强收敛蕴涵着依范数收敛. 事实上, 这可从三角不等式的以下推论得到:
\begin{align*}
\text{对所有 }X\text{ 中的 }g,h,\quad|\|g\|-\|h\||\leqslant\|g-h\|.
\end{align*}
剩下来要证明在 $L^{p}(E)$ ($1<p<\infty$) 空间中, 弱收敛和依范数收敛蕴涵在 $L^{p}(E)$ 中强收敛. 我们对 $p=2$ 的情形给出证明\footnote{对一般 $p>1$ 的证明, 等式 $(a-b)^{2}=a^{2}-2ab+b^{2}$ 需要替换. 详细的证明在 Frigyes Riesz 和 Bela Sz.-Nagy 的 \emph{Functional Analysis}[RSN90] 78$\sim$80 页给出.}. 令 $\{f_{n}\}$ 为 $L^{2}(E)$ 中的序列, 使得
\begin{align*}
\text{在 }L^{2}(E)\text{ 中 }\{f_{n}\}\rightharpoonup f\text{ 且 }\lim_{n\to\infty}\int_{E}f_{n}^{2}=\int_{E}f^{2}.
\end{align*}
观察到对每个 $n$,
\begin{align*}
\|f_{n}-f\|_{2}^{2}=\int_{E}|f_{n}-f|^{2}=\int_{E}(f_{n}-f)^{2}=\int_{E}|f_{n}|^{2}-2\cdot\int_{E}f_{n}\cdot f+\int_{E}|f|^{2}.
\end{align*}
由于 $f$ 属于 $L^{q}(E)=L^{2}(E)$,
\begin{align*}
\lim_{n\to\infty}\int_{E}f_{n}\cdot f=\int_{E}f^{2}.
\end{align*}
因此在 $L^{2}(E)$ 中 $\{f_{n}\}\to f$.
\end{proof}

\begin{corollary}\label{corollary:ch08-弱收敛子列强收敛}
令 $E$ 为可测集而 $1<p<\infty$. 假定在 $L^{p}(E)$ 中 $\{f_{n}\}\rightharpoonup f$. 则存在 $\{f_{n}\}$ 的子序列在 $L^{p}(E)$ 中强收敛于 $f$ 当且仅当
\begin{align*}
\|f\|_{p}=\liminf\|f_{n}\|_{p}.
\end{align*}
\end{corollary}

\begin{proof}
若 $\|f\|_{p}=\liminf\|f_{n}\|_{p}$, 则存在子序列 $\{f_{n_{k}}\}$ 使得 $\lim_{k\to\infty}\|f_{n_{k}}\|_{p}=\|f\|_{p}$. Radon-Riesz 定理告诉我们 $\{f_{n_{k}}\}$ 在 $L^{q}(E)$ 中强收敛于 $f$. 反过来, 若存在强收敛于 $f$ 的子序列 $\{f_{n_{k}}\}$, 则 $\lim_{k\to\infty}\|f_{n_{k}}\|_{p}=\|f\|_{p}$. 因此 $\liminf\|f_{n}\|_{p}\leqslant\|f\|_{p}$. \eqref{align:ch08-弱收敛有界}右边的不等式是该不等式的反向不等式.
\end{proof}

如同下面例子所表明的, Radon-Riesz 定理不能推广到 $p=1$ 的情形.

\begin{example}\label{example:ch08-Radon反例}
对每个自然数 $n$, 在 $I=[-\pi,\pi]$ 上定义 $f_{n}(x)=1+\sin(nx)$. 从 Riemann-Lebesgue 引理得出序列 $\{f_{n}\}$ 在 $L^{1}(I)$ 中弱收敛于函数 $f\equiv 1$. 由于每个函数 $f_{n}$ 非负, 因此有 $\lim_{n\to\infty}\|f_{n}\|_{1}=\|f\|_{1}$. 由于 $\{\sin(nx)\}$ 在 $L^{1}(I)$ 中不强收敛于 $f\equiv 0$, $\{f_{n}\}$ 在 $L^{1}(I)$ 中不强收敛于 $f\equiv 1$.
\end{example}

\begin{remark}\label{remark:ch08-共轭函数泛函}
令 $E$ 为可测集, $1\leqslant p<\infty$, $f\in L^{p}(E)$, $f\neq 0$, 而 $f^{*}$ 为 $f$ 的共轭函数, 定义 $T\in(L^{p}(E))^{*}$ 为
\begin{align}
\text{对所有 }h\in L^{q}(E),\quad T(h)=\int_{E}f^{*}\cdot h.\label{align:ch08-共轭函数泛函}
\end{align}
根据上一章的定理 1 和本章的命题 2,
\begin{align}
T(f)=\|f\|_{p}\text{ 且 }\|T\|_{*}=1.\label{align:ch08-共轭函数泛函范数}
\end{align}
在 14.2 节, 我们证明 Hahn-Banach 定理, 而作为该定理的一个推论, 它表明了若 $X$ 是任何赋范线性空间而 $f$ 属于 $X$, 则存在 $X^{*}$ 中的有界线性算子 $T$ 使得 $T(f)=\|f\|$ 且 $\|T\|_{*}=1$. 对于 $L^{p}(E)$ 空间, 对共轭函数的积分给出了该抽象泛函的一个具体表示.
\end{remark}

\begin{exercise}\label{exercise:ch08-12}
证明本节第一个例子中定义的序列对所有 $1\leqslant p<\infty$ 在 $L^{p}[0,1]$ 中不强收敛于 $f\equiv 0$.
\end{exercise}

\begin{exercise}\label{exercise:ch08-13}
固定实数 $\alpha$ 和 $\beta$. 对每个自然数 $k$ 和 $n$, 考虑定义在 $I=[0,1]$ 上的阶梯函数 $\{f_{n}\}$:
\begin{align*}
f_{n}(x)=(1-(-1)^{k})\alpha/2+(1+(-1)^{k})\beta/2,\text{ 对于 }k/2^{n}\leqslant x<(k+1)/2^{n},\,0\leqslant k<2^{n}-1.
\end{align*}
对 $1<p<\infty$, 证明 $\{f_{n}\}$ 在 $L^{p}(I)$ 中弱收敛到取值 $(\alpha+\beta)/2$ 的常函数. 对于 $\alpha\neq\beta$, 证明 $\{f_{n}\}$ 没有在 $L^{p}(I)$ 中强收敛的子序列.
\end{exercise}

\begin{exercise}\label{exercise:ch08-14}
令 $h$ 为定义在整个 $\mathbf{R}$ 上的连续函数, 它是以 $T$ 为周期的周期函数且 $\int_{0}^{T}h=0$. 令 $[a,b]$ 为闭有界区间且对每个自然数 $n$, 定义 $[a,b]$ 上的函数 $f_{n}$ 为 $f_{n}(x)=h(nx)$. 在 $[a,b]$ 上定义 $f\equiv 0$. 证明对于 $1\leqslant p<\infty$, 在 $L^{p}[a,b]$ 中 $\{f_{n}\}\rightharpoonup f$.
\end{exercise}

\begin{exercise}\label{exercise:ch08-15}
令 $1<p<\infty$ 而 $f_{0}$ 属于 $L^{p}(\mathbf{R})$. 对每个自然数 $n$, 对所有 $x$ 定义 $f_{n}(x)=f_{0}(x-n)$. 在 $\mathbf{R}$ 上定义 $f\equiv 0$. 证明在 $L^{p}(\mathbf{R})$ 中 $\{f_{n}\}\rightharpoonup f$. 这对 $p=1$ 成立吗?
\end{exercise}

\begin{exercise}\label{exercise:ch08-16}
令 $E$ 为可测集, $\{f_{n}\}$ 是 $L^{2}(E)$ 中的序列而 $f$ 属于 $L^{2}(E)$. 假定
\begin{align*}
\lim_{n\to\infty}\int_{E}f_{n}\cdot f=\lim_{n\to\infty}\int_{E}f_{n}^{2}=\int_{E}f^{2}.
\end{align*}
证明在 $L^{2}(E)$ 中 $\{f_{n}\}$ 强收敛于 $f$.
\end{exercise}

\begin{exercise}\label{exercise:ch08-17}
令 $E$ 为可测集而 $1<p<\infty$. 假定 $\{f_{n}\}$ 是 $L^{p}(E)$ 中的有界序列而 $f$ 属于 $L^{p}(E)$. 考虑以下四个性质: (i) $\{f_{n}\}$ 在 $E$ 上几乎处处逐点收敛于 $f$, (ii) 在 $L^{p}(E)$ 中 $\{f_{n}\}\rightharpoonup f$, (iii) $\{\|f_{n}\|_{p}\}$ 收敛于 $\|f\|_{p}$, (iv) 在 $L^{p}(E)$ 中 $\{f_{n}\}\to f$. 若 $\{f_{n}\}$ 具有这些性质中的两个, 是否存在子序列具有所有四个性质?
\end{exercise}

\begin{exercise}\label{exercise:ch08-18}
令 $X$ 为赋范线性空间而在 $X$ 中 $\{f_{n}\}\rightharpoonup f$. 假定 $\{\|f_{n}\|\}$ 是无界的. 证明通过可能地取子序列以及重新标记, 可以假定对所有 $n$, $\|f_{n}\|\geqslant a_{n}=n\cdot 3^{n}$. 接着证明通过可能地取进一步的子序列以及重新标记, 可以假定 $\{\|f_{n}\|/a_{n}\}\to\alpha\in[1,\infty]$. 对每个 $n$, 定义 $g_{n}=\alpha_{n}/\|f_{n}\|\cdot f_{n}$. 证明 $\{g_{n}\}$ 弱收敛且对所有 $n$, $\|g_{n}\|=n\cdot 3^{n}$.
\end{exercise}

\begin{exercise}\label{exercise:ch08-19}
对 $1\leqslant p<\infty$, 令 $\{\xi_{n}\}$ 为 $\ell^{p}$ 中的有界序列且 $\xi$ 属于 $\ell^{p}$. 证明在 $\ell^{p}$ 中 $\{\xi_{n}\}\rightharpoonup\xi$ 当且仅当它依分量收敛, 即
\begin{align*}
\text{对每个指标 }k,\quad\lim_{n\to\infty}\xi_{n}^{k}=\xi^{k},\quad\text{其中 }\xi_{n}=<\xi_{n}^{1},\xi_{n}^{2},\cdots>,\quad\text{而 }\xi=<\xi^{1},\xi^{2},\cdots>.
\end{align*}
\end{exercise}

\begin{exercise}\label{exercise:ch08-20}
令 $1\leqslant p_{1}<p_{2}<\infty$, $\{f_{n}\}$ 为 $L^{p_{2}}[0,1]$ 中的序列而 $f$ 属于 $L^{p_{2}}[0,1]$. 在 $L^{p_{2}}[0,1]$ 中 $\{f_{n}\}\rightharpoonup f$ 与在 $L^{p_{1}}[0,1]$ 中 $\{f_{n}\}\rightharpoonup f$ 的关系是什么?
\end{exercise}

\begin{exercise}\label{exercise:ch08-21}
对于 $1\leqslant p<\infty$ 和每个指标 $n$, 令 $e_{n}\in\ell^{p}$ 的第 $n$ 个分量为 $1$, 而其他分量等于零. 证明: 若 $p>1$, 则 $\{e_{n}\}$ 在 $\ell^{p}$ 中弱收敛于 $0$, 但没有强收敛到 $0$ 的子序列. 证明 $\{e_{n}\}$ 在 $\ell^{1}$ 中不弱收敛.
\end{exercise}

\begin{exercise}\label{exercise:ch08-22}
在 $\ell^{2}$ 中叙述与证明 Radon-Riesz 定理.
\end{exercise}

\begin{exercise}\label{exercise:ch08-23}
令 $[a,b]$ 为闭有界区间. 假定在 $C[a,b]$ 中 $\{f_{n}\}\rightharpoonup f$. 证明 $\{f_{n}\}$ 在 $[a,b]$ 上逐点收敛于 $f$.
\end{exercise}

\begin{exercise}\label{exercise:ch08-24}
令 $[a,b]$ 为闭有界区间. 假定在 $L^{\infty}[a,b]$ 中 $\{f_{n}\}\rightharpoonup f$. 证明对所有 $x\in[a,b]$,
\begin{align*}
\lim_{n\to\infty}\int_{a}^{x}f_{n}=\int_{a}^{x}f.
\end{align*}
\end{exercise}

\begin{exercise}\label{exercise:ch08-25}
令 $X$ 为赋范线性空间. 假定对每个 $f\in X$, 存在有界线性泛函 $T\in X^{*}$, 使得 $T(f)=\|f\|$ 且 $\|T\|_{*}=1$.
\begin{enumerate}[(1)]
\item 证明: 若 $\{f_{n}\}$ 在 $X$ 中弱收敛于 $f_{1}$ 与 $f_{2}$, 则 $f_{1}=f_{2}$.
\item 证明: 若在 $X$ 中 $\{f_{n}\}\rightharpoonup f$, 则 $\|f\|\leqslant\liminf\|f_{n}\|$.
\end{enumerate}
\end{exercise}

\begin{exercise}\label{exercise:ch08-26}
(一致有界原理) 令 $E$ 为可测集, $1\leqslant p<\infty$, 而 $q$ 是 $p$ 的共轭. 假定 $\{f_{n}\}$ 是 $L^{p}(E)$ 中的序列使得对每个 $g\in L^{q}(E)$, 序列 $\{\int_{E}g\cdot f_{n}\}$ 是有界的, 证明 $\{f_{n}\}$ 在 $L^{p}(E)$ 中是有界的.
\end{exercise}

\end{document}